% MJ46--1C
% Akuter Asthmaanfall
% 14.0
% 2013-11-30

:ref:`m-asthmaanfall`
Taktik Linderung der Atemnot und rasche
medikamentöse Therapie **Vitale Bedrohung im Falle
bleibender schwerer Atemnot bzw. Status asthmaticus**
}

\begin{itemize}
  -   Lagerung: **Oberkörper hoch**, abstützen lassen
  (Unterstützung der Atemhilfsmuskulatur)
  -   O₂-Gabe gemäß
  \prettyref{m:sauerstoffberieselung} 
  
  {Achtung: In seltenen Fällen
  kann es vorkommen, dass bei hochdosierter
  Sauerstoffgabe der Patient weniger atmet und
  Bewusstseinsstörungen auftreten (CO₂-Narkose)! Siehe
  \enquote{\Abk{COPD}} (\FullPrettyRef{topic:copd}). Bei sorgfältiger
  Überwachung des Patienten stellt die O₂-Gabe in der
  Praxis jedoch kein Problem
  dar \citep{ScholzNotfallmedizin2}.}
  \hfill\Commentary[Asthma, Sauerstoffgabe]{Asthma:
  \protectO₂ bis eine Sauerstoffsättigung von \VonBis{94}{98}\PC* erreicht ist, oder 8 L / min, gem.
 :cite:`AsboeLehrmeinungRd2011-11`\ . Die einschlägige
  Literatur ist widersprüchlich. Besonderes betont muss die Notwendigkeit zur Überwachung des
  Patienten bezüglich Bewusstseinsstörungen und Atmung
  werden!)}
  -   Beruhigen, voratmen, durch fast
  geschlossene Lippen ausatmen lassen (Lippenbremse)
  -   Sprays, die Patient evtl. bei sich hat, dürfen
  nicht mehr genommen werden.
  \hfill\Commentary[Asthma / Spray]{Oft stellt sich die
  Frage \enquote{Darf der Patient seinen (vom Arzt vielleicht gerade für
  diese Situation verschriebenen) Spray nehmen?}
  Grundsätzlich soll beim schweren Asthma-Anfall davon
  abgesehen werden, \enquote{blind} ein inhalatives Beta-Mimetikum
  zu geben, das zweifelhaft ist, ob die Luftpassage
  ausreicht, umd den Wirkstoff in die Lunge zu befördern.
  Eine generelle Freigabe, verschriebene Srpays einzunehmen,
  kann in Anbetracht der Nebenwirkungen bei
  Schleimhautresorption nicht gegeben werden.}
\end{itemize}

\Customized{
\begin{description}
  \item[\CustAsboe] Beim akuten Bronchospasmus ist die Gabe
  von eines Beta-Mimetikums gem. Algorithmus durch NFS vorgesehen.
 :cite:`Asb921Memo-AL-2011-000001`\ .
\end{description}
}